%=====================================================================
% RH-NBV planner — lean version.
% Keeps the original single-algorithm structure; adds only what the
% contribution claims require: the redundancy penalty (2 lines), a
% one-line stagnation-escape summary, and workspace projection (1 line).
% Full escape details belong in the text (see the \Comment references).
% Requires: \usepackage{algorithm} \usepackage{algpseudocode} \usepackage{amsmath}
%=====================================================================

\begin{algorithm}[t]
\caption{Receding-Horizon Next-Best-View Planning}
\label{alg:rh-nbv}
\begin{algorithmic}[1]
\State \textbf{Input:} current configuration $\xi_0$, current map $M_0$
\State \textbf{Parameters:} horizon $H$, candidates $K$, discount $\gamma$,
       cost weight $\lambda$, target $\mathbf{t}$, step size $s$,
       orbit bias $b$, camera workspace $\mathcal{B}$
\Statex
\If{ROI coverage has stagnated for $P$ iterations}
    \State relocate $\xi_0$ to a distant random position in $\mathcal{B}$
           \Comment{escape, Sec.~\ref{sec:stagnation}}
\EndIf
\State $J_{\mathrm{best}} \gets -\infty$;\quad $\Xi_{\mathrm{best}} \gets \emptyset$
\For{$i = 0$ \textbf{to} $K-1$}
    \State $\Xi \gets$ \Call{GenerateCandidateSequence}{$\xi_0, H$}
    \State $J \gets$ \Call{EvaluateSequence}{$\Xi, \xi_0, M_0$}
    \If{$J > J_{\mathrm{best}}$}
        \State $J_{\mathrm{best}} \gets J$;\quad $\Xi_{\mathrm{best}} \gets \Xi$
    \EndIf
\EndFor
\State $\xi_{\mathrm{next}} \gets \Xi_{\mathrm{best}}[0]$
       \Comment{execute first step only, then replan}
\State $\mathbf{q}_{\mathrm{next}} \gets$ \Call{LookAt}{$\xi_{\mathrm{next}}, \mathbf{t}$}
\State \Return $(\xi_{\mathrm{next}}, \mathbf{q}_{\mathrm{next}})$
\Statex
\Procedure{GenerateCandidateSequence}{$\xi_0, H$}
    \State $\mathrm{prev} \gets \xi_0$
    \For{$k = 0$ \textbf{to} $H-1$}
        \If{$\mathrm{rand}() < b$}
            \State $\mathbf{r} \gets \mathrm{normalize}(\mathbf{t} - \mathrm{prev})$
            \State $\mathbf{v} \gets$ random unit vector perpendicular to $\mathbf{r}$
            \State $\xi_k \gets \mathrm{prev} + \mathbf{v} \cdot
                   \bigl(s \cdot \mathrm{uniform}(0.3, 1.0)\bigr)$
        \Else
            \State $\xi_k \gets \mathrm{uniform\_sample}(\mathcal{B})$
        \EndIf
        \State project $\xi_k$ to the reachable workspace (min.\ stand-off, reach bounds)
        \State $\Xi[k] \gets \xi_k$;\quad $\mathrm{prev} \gets \xi_k$
    \EndFor
    \State \Return $\Xi$
\EndProcedure
\Statex
\Procedure{EvaluateSequence}{$\Xi, \xi_0, M_0$}
    \State $J \gets 0$;\quad $\mathrm{prev} \gets \xi_0$;\quad
           $M_{\mathrm{pred}} \gets \mathrm{copy}(M_0)$
    \For{$k = 0$ \textbf{to} $H-1$}
        \State $f \gets$ \Call{ComputeGain}{$M_{\mathrm{pred}}, \Xi[k]$}
        \If{$\min_{j<k} \lVert \Xi[k] - \Xi[j] \rVert_2 < s/2$}
            \State $f \gets 0.3\,f$ \Comment{redundancy penalty}
        \EndIf
        \State $J \gets J + \gamma^{k}\,
               \bigl(f - \lambda \lVert \Xi[k] - \mathrm{prev} \rVert_2\bigr)$
        \State $M_{\mathrm{pred}} \gets$
               \Call{PredictUpdate}{$M_{\mathrm{pred}}, \Xi[k]$}
               \Comment{belief forward simulation}
        \State $\mathrm{prev} \gets \Xi[k]$
    \EndFor
    \State \Return $J$
\EndProcedure
\end{algorithmic}
\end{algorithm}

%---------------------------------------------------------------------
\subsection{Algorithm Walkthrough}
\label{sec:alg-walkthrough}

Each planning iteration of Algorithm~\ref{alg:rh-nbv} proceeds in three
stages. First, a \emph{stagnation guard} checks whether ROI coverage has
plateaued; if so, the camera is relocated to a distant reachable position
before any planning takes place, preventing the local sampler from
repeatedly drawing candidates inside an exhausted region
(Sec.~\ref{sec:stagnation}). Second, $K$ candidate view sequences of
length $H$ are generated by \textsc{GenerateCandidateSequence}: with
probability $b$ a step moves tangentially on the sphere around the target
$\mathbf{t}$---the component of camera motion that changes the viewing
angle rather than merely the viewing distance---and with probability
$1-b$ it jumps to a uniformly random position in the workspace
$\mathcal{B}$, injecting the diversity needed to escape locally optimal
orbits. Every step is projected back to the reachable workspace so that
no candidate violates the minimum stand-off distance or the robot reach.
Third, \textsc{EvaluateSequence} scores each sequence with the discounted
objective of Eq.~\ref{eq:rw-ours}: the semantic information gain of every
step is computed on a \emph{predicted} map $M_{\mathrm{pred}}$ that is
rolled forward after each hypothetical view, so that later steps are only
credited for information not already gathered by earlier ones, and views
closer than $s/2$ to a previous step of the same sequence incur a
redundancy penalty. Finally, only the first viewpoint of the best
sequence is executed; the remaining $H-1$ steps are discarded and
replanned after the next real measurement, which is what makes the
planner robust to prediction errors in $M_{\mathrm{pred}}$.

%---------------------------------------------------------------------
\subsection{Parameter Settings}
\label{sec:alg-parameters}

Table~\ref{tab:rh-params} lists all planner parameters, the values used
in the experiments, and the rationale behind each choice.

\begin{table}[t]
  \centering
  \caption{RH--NBV planner parameters used in all experiments.}
  \label{tab:rh-params}
  \begin{tabular}{clcl}
    \toprule
    Symbol & Meaning & Value & Rationale \\
    \midrule
    $H$ & planning horizon & $3$ &
      ablation over $H \in \{1,2,3,4\}$; best trade-off \\
    $K$ & candidate sequences & $10$ &
      matches the candidate count of SamplingNBV \\
    $\gamma$ & discount factor & $0.85$ &
      later steps rely on predicted maps; trust less \\
    $\lambda$ & motion-cost weight & $2.0$ &
      penalises workspace-crossing moves \\
    $s$ & step size & $0.065$\,m &
      identical to GradientNBV's step size \\
    $b$ & orbit bias & $0.7$ &
      exploitation/exploration balance \\
    $P$ & stagnation patience & $4$ iters &
      tolerates noise; caps wasted budget \\
    $\tau$ & stagnation threshold & $1.5\,\%$ &
      minimum coverage gain counted as progress \\
    \bottomrule
  \end{tabular}
\end{table}

The two RH-specific extensions---the spherical orbit shell and the
occlusion-bonus term---are \emph{disabled} in all baseline comparisons,
since neither has a counterpart in GradientNBV; they are evaluated
separately as ablations. All stochastic components use a fixed random
seed so that every experiment is exactly reproducible.

%---------------------------------------------------------------------
\subsection{Computational Cost}
\label{sec:alg-complexity}

The dominant cost of one planning iteration is ray-traced gain
evaluation. Algorithm~\ref{alg:rh-nbv} performs exactly $K \cdot H$ calls
to \textsc{ComputeGain} ($30$ with the settings of
Table~\ref{tab:rh-params}), each casting one virtual image
($600 \times 450$ rays, $128$ samples per ray) through the voxel grid on
the GPU. In addition, $K \cdot H$ calls to \textsc{PredictUpdate} perform
a ray traversal of comparable size; following the cost metric of
\textcite{burusa2024}, which counts information-gain evaluations only,
these are not included in the reported ray-tracing-call counts, but their
wall-clock cost is included in all timing results. Per iteration, the
planner is therefore roughly $K \cdot H$ times more expensive than
GradientNBV (one utility evaluation per gradient step) and $\sim\!3\times$
more expensive than SamplingNBV with $10$ candidates ($10$ evaluations);
the experiments in Chapter~\ref{ch:experimental-setup} quantify whether
the non-myopic sequence evaluation justifies this overhead. All gain and
prediction computations run under \texttt{torch.no\_grad()}, since the
sampling-based search requires no gradients; this halves memory traffic
relative to a na\"ive autograd implementation without changing any
numerical result.
